\documentclass[12pt]{article}
\usepackage[UTF8]{inputenc}
\usepackage{xeCJK}
\setCJKmainfont{Sarabun}
\usepackage{enumitem}
\usepackage{geometry}
\geometry{a4paper, margin=2.5cm}
\usepackage{setspace}
\onehalfspacing

\begin{document}

\begin{center}
{\Large\textbf{กลศาสตร์ควอนตัมเบื้องต้น}}\\[6pt]
{\normalsize แบบฝึกหัดเพิ่มเติม}\\[6pt]
{\normalsize วิชาฟิสิกส์ ม.ปลาย บทที่ 26}\\[4pt]
{\footnotesize ทั้งหมด 30 ข้อ}
\end{center}
\hrule
\bigskip

\section*{แบบฝึกหัดเพิ่มเติม}
\begin{enumerate}[label=\arabic*., itemsep=0.5\baselineskip, topsep=0.5\baselineskip, leftmargin=*]
\item de Broglie ได้รับรางวัลโนเบลสาขาฟิสิกส์ปีใด?
\begin{enumerate}[label=\alph*)., itemsep=0.3\baselineskip, topsep=0.3\baselineskip]
\item ก\ 
\begin{minipage}[t]{0.9\linewidth}
\raggedright 1918
\end{minipage}
\item ข\ 
\begin{minipage}[t]{0.9\linewidth}
\raggedright 1929
\end{minipage}
\item ค\ 
\begin{minipage}[t]{0.9\linewidth}
\raggedright 1933
\end{minipage}
\item ง\ 
\begin{minipage}[t]{0.9\linewidth}
\raggedright 1945
\end{minipage}
\end{enumerate}
\vspace{-0.3\baselineskip}
\begin{small}
\textit{คำตอบ: ข}
\end{small}

\item ใครอธิบายปรากฏการณ์โฟโตอิเล็กทริกได้อย่างถูกต้องและได้รับรางวัลโนเบลจากผลงานนี้?
\begin{enumerate}[label=\alph*)., itemsep=0.3\baselineskip, topsep=0.3\baselineskip]
\item ก\ 
\begin{minipage}[t]{0.9\linewidth}
\raggedright Planck
\end{minipage}
\item ข\ 
\begin{minipage}[t]{0.9\linewidth}
\raggedright Einstein
\end{minipage}
\item ค\ 
\begin{minipage}[t]{0.9\linewidth}
\raggedright Bohr
\end{minipage}
\item ง\ 
\begin{minipage}[t]{0.9\linewidth}
\raggedright Heisenberg
\end{minipage}
\end{enumerate}
\vspace{-0.3\baselineskip}
\begin{small}
\textit{คำตอบ: ข}
\end{small}

\item พลังงานของโฟตอนแสงสีเขียว (λ = 550 nm) มีค่าเท่าใด?
\begin{enumerate}[label=\alph*)., itemsep=0.3\baselineskip, topsep=0.3\baselineskip]
\item ก\ 
\begin{minipage}[t]{0.9\linewidth}
\raggedright 1.8 eV
\end{minipage}
\item ข\ 
\begin{minipage}[t]{0.9\linewidth}
\raggedright 2.3 eV
\end{minipage}
\item ค\ 
\begin{minipage}[t]{0.9\linewidth}
\raggedright 3.6 eV
\end{minipage}
\item ง\ 
\begin{minipage}[t]{0.9\linewidth}
\raggedright 4.5 eV
\end{minipage}
\end{enumerate}
\vspace{-0.3\baselineskip}
\begin{small}
\textit{คำตอบ: ข}
\end{small}

\item ไฮเซนเบิร์กได้รับรางวัลโนเบลสาขาฟิสิกส์ปีใด?
\begin{enumerate}[label=\alph*)., itemsep=0.3\baselineskip, topsep=0.3\baselineskip]
\item ก\ 
\begin{minipage}[t]{0.9\linewidth}
\raggedright 1918
\end{minipage}
\item ข\ 
\begin{minipage}[t]{0.9\linewidth}
\raggedright 1932
\end{minipage}
\item ค\ 
\begin{minipage}[t]{0.9\linewidth}
\raggedright 1945
\end{minipage}
\item ง\ 
\begin{minipage}[t]{0.9\linewidth}
\raggedright 1958
\end{minipage}
\end{enumerate}
\vspace{-0.3\baselineskip}
\begin{small}
\textit{คำตอบ: ข}
\end{small}

\item ใครเสนอว่าฟังก์ชันคลื่นบอก 'ความน่าจะเป็น' ในการพบอนุภาค?
\begin{enumerate}[label=\alph*)., itemsep=0.3\baselineskip, topsep=0.3\baselineskip]
\item ก\ 
\begin{minipage}[t]{0.9\linewidth}
\raggedright de Broglie
\end{minipage}
\item ข\ 
\begin{minipage}[t]{0.9\linewidth}
\raggedright Schrödinger
\end{minipage}
\item ค\ 
\begin{minipage}[t]{0.9\linewidth}
\raggedright Born
\end{minipage}
\item ง\ 
\begin{minipage}[t]{0.9\linewidth}
\raggedright Heisenberg
\end{minipage}
\end{enumerate}
\vspace{-0.3\baselineskip}
\begin{small}
\textit{คำตอบ: ค}
\end{small}

\item ถ้าอิเล็กตรออนเคลื่อนที่ช้าลง ความยาวคลื่น de Broglie จะเป็นอย่างไร?
\begin{enumerate}[label=\alph*)., itemsep=0.3\baselineskip, topsep=0.3\baselineskip]
\item ก\ 
\begin{minipage}[t]{0.9\linewidth}
\raggedright เพิ่มขึ้น
\end{minipage}
\item ข\ 
\begin{minipage}[t]{0.9\linewidth}
\raggedright ลดลง
\end{minipage}
\item ค\ 
\begin{minipage}[t]{0.9\linewidth}
\raggedright คงที่
\end{minipage}
\item ง\ 
\begin{minipage}[t]{0.9\linewidth}
\raggedright เป็นศูนย์
\end{minipage}
\end{enumerate}
\vspace{-0.3\baselineskip}
\begin{small}
\textit{คำตอบ: ก}
\end{small}

\item ค่างาน (work function) ของโลหะขึ้นอยู่กับอะไร?
\begin{enumerate}[label=\alph*)., itemsep=0.3\baselineskip, topsep=0.3\baselineskip]
\item ก\ 
\begin{minipage}[t]{0.9\linewidth}
\raggedright ความถี่ของแสงที่ตกกระทบ
\end{minipage}
\item ข\ 
\begin{minipage}[t]{0.9\linewidth}
\raggedright ชนิดของโลหะ
\end{minipage}
\item ค\ 
\begin{minipage}[t]{0.9\linewidth}
\raggedright ความเข้มของแสง
\end{minipage}
\item ง\ 
\begin{minipage}[t]{0.9\linewidth}
\raggedright อุณหภูมิของแสง
\end{minipage}
\end{enumerate}
\vspace{-0.3\baselineskip}
\begin{small}
\textit{คำตอบ: ข}
\end{small}

\item แสงที่มีความยาวคลื่น 500 nm มีความถี่เท่าใด? (c = 3×10⁸ m/s)
\begin{enumerate}[label=\alph*)., itemsep=0.3\baselineskip, topsep=0.3\baselineskip]
\item ก\ 
\begin{minipage}[t]{0.9\linewidth}
\raggedright 5 × 10¹⁴ Hz
\end{minipage}
\item ข\ 
\begin{minipage}[t]{0.9\linewidth}
\raggedright 6 × 10¹⁴ Hz
\end{minipage}
\item ค\ 
\begin{minipage}[t]{0.9\linewidth}
\raggedright 3 × 10¹⁴ Hz
\end{minipage}
\item ง\ 
\begin{minipage}[t]{0.9\linewidth}
\raggedright 1.5 × 10¹⁵ Hz
\end{minipage}
\end{enumerate}
\vspace{-0.3\baselineskip}
\begin{small}
\textit{คำตอบ: ข}
\end{small}

\item สมการ Schrödinger แบบไม่ขึ้นกับเวลา (time-independent Schrödinger equation) ใช้อธิบายอะไร?
\begin{enumerate}[label=\alph*)., itemsep=0.3\baselineskip, topsep=0.3\baselineskip]
\item ก\ 
\begin{minipage}[t]{0.9\linewidth}
\raggedright พลังงานของคลื่นแม่เหล็กไฟฟ้า
\end{minipage}
\item ข\ 
\begin{minipage}[t]{0.9\linewidth}
\raggedright สถานะคงที่ (stationary states) ของอะตอม
\end{minipage}
\item ค\ 
\begin{minipage}[t]{0.9\linewidth}
\raggedright การเคลื่อนที่ของดาวเคราะห์
\end{minipage}
\item ง\ 
\begin{minipage}[t]{0.9\linewidth}
\raggedright แรงโนบลังก์
\end{minipage}
\end{enumerate}
\vspace{-0.3\baselineskip}
\begin{small}
\textit{คำตอบ: ข}
\end{small}

\item การกระเจิงของอิเล็กตรออนโดยผลึก (electron diffraction) เป็นปรากฏการณ์ที่ยืนยันอะไร?
\begin{enumerate}[label=\alph*)., itemsep=0.3\baselineskip, topsep=0.3\baselineskip]
\item ก\ 
\begin{minipage}[t]{0.9\linewidth}
\raggedright อิเล็กตรออนเป็นอนุภาค
\end{minipage}
\item ข\ 
\begin{minipage}[t]{0.9\linewidth}
\raggedright อิเล็กตรออนมีสมบัติคลื่น
\end{minipage}
\item ค\ 
\begin{minipage}[t]{0.9\linewidth}
\raggedright อิเล็กตรออนเป็นทรงกลม
\end{minipage}
\item ง\ 
\begin{minipage}[t]{0.9\linewidth}
\raggedright อิเล็กตรออนมีประจุ
\end{minipage}
\end{enumerate}
\vspace{-0.3\baselineskip}
\begin{small}
\textit{คำตอบ: ข}
\end{small}

\item ถ้าแสงความถี่ต่ำกว่าความถี่ขีดเริ่มตกกระทบโลหะ จะเกิดอะไรขึ้น?
\begin{enumerate}[label=\alph*)., itemsep=0.3\baselineskip, topsep=0.3\baselineskip]
\item ก\ 
\begin{minipage}[t]{0.9\linewidth}
\raggedright อิเล็กตรออนหลุดออกมาทันที
\end{minipage}
\item ข\ 
\begin{minipage}[t]{0.9\linewidth}
\raggedright อิเล็กตรออนหลุดออกมาหลังผ่านไปสักครู่
\end{minipage}
\item ค\ 
\begin{minipage}[t]{0.9\linewidth}
\raggedright ไม่มีอิเล็กตรออนหลุดไม่ว่าความเข้มจะเท่าใด
\end{minipage}
\item ง\ 
\begin{minipage}[t]{0.9\linewidth}
\raggedright อิเล็กตรออนหลุดน้อยลง
\end{minipage}
\end{enumerate}
\vspace{-0.3\baselineskip}
\begin{small}
\textit{คำตอบ: ค}
\end{small}

\item แสงสีใดมีโฟตอนพลังงานสูงที่สุดระหว่าง แดง เขียว น้ำเงิน เหลือง?
\begin{enumerate}[label=\alph*)., itemsep=0.3\baselineskip, topsep=0.3\baselineskip]
\item ก\ 
\begin{minipage}[t]{0.9\linewidth}
\raggedright แดง
\end{minipage}
\item ข\ 
\begin{minipage}[t]{0.9\linewidth}
\raggedright เขียว
\end{minipage}
\item ค\ 
\begin{minipage}[t]{0.9\linewidth}
\raggedright น้ำเงิน
\end{minipage}
\item ง\ 
\begin{minipage}[t]{0.9\linewidth}
\raggedright เหลือง
\end{minipage}
\end{enumerate}
\vspace{-0.3\baselineskip}
\begin{small}
\textit{คำตอบ: ค}
\end{small}

\item ข้อใดเป็นความสัมพันธ์ที่ถูกต้องของหลักความไม่แน่นอน?
\begin{enumerate}[label=\alph*)., itemsep=0.3\baselineskip, topsep=0.3\baselineskip]
\item ก\ 
\begin{minipage}[t]{0.9\linewidth}
\raggedright Δx · Δp = ħ/2
\end{minipage}
\item ข\ 
\begin{minipage}[t]{0.9\linewidth}
\raggedright Δx · Δp ≥ ħ/2
\end{minipage}
\item ค\ 
\begin{minipage}[t]{0.9\linewidth}
\raggedright Δx · Δp ≤ ħ/2
\end{minipage}
\item ง\ 
\begin{minipage}[t]{0.9\linewidth}
\raggedright Δx · Δp > ħ/2 เสมอ
\end{minipage}
\end{enumerate}
\vspace{-0.3\baselineskip}
\begin{small}
\textit{คำตอบ: ข}
\end{small}

\item ทำไมกล้องจุลทรรศน์อิเล็กตรออน (Electron Microscope) จึงมีกำลังขยายสูงกว่ากล้องจุลทรรศน์แสง?
\begin{enumerate}[label=\alph*)., itemsep=0.3\baselineskip, topsep=0.3\baselineskip]
\item ก\ 
\begin{minipage}[t]{0.9\linewidth}
\raggedright เพราะอิเล็กตรออนเล็กกว่าแสง
\end{minipage}
\item ข\ 
\begin{minipage}[t]{0.9\linewidth}
\raggedright เพราะความยาวคลื่นของอิเล็กตรออนสั้นกว่าแสงที่ตามองเห็น ทำให้รายละเอียดเล็กๆ ถูกเลี้ยวเบนได้น้อยกว่า
\end{minipage}
\item ค\ 
\begin{minipage}[t]{0.9\linewidth}
\raggedright เพราะอิเล็กตรออนเคลื่อนที่เร็วกว่าแสง
\end{minipage}
\item ง\ 
\begin{minipage}[t]{0.9\linewidth}
\raggedright เพราะอิเล็กตรออนมีมวล
\end{minipage}
\end{enumerate}
\vspace{-0.3\baselineskip}
\begin{small}
\textit{คำตอบ: ข}
\end{small}

\item อิเล็กตรออนที่มีพลังงานจลน์ 150 eV มีความยาวคลื่น de Broglie เท่าใด?
\begin{enumerate}[label=\alph*)., itemsep=0.3\baselineskip, topsep=0.3\baselineskip]
\item ก\ 
\begin{minipage}[t]{0.9\linewidth}
\raggedright 0.01 nm
\end{minipage}
\item ข\ 
\begin{minipage}[t]{0.9\linewidth}
\raggedright 0.1 nm
\end{minipage}
\item ค\ 
\begin{minipage}[t]{0.9\linewidth}
\raggedright 1 nm
\end{minipage}
\item ง\ 
\begin{minipage}[t]{0.9\linewidth}
\raggedright 10 nm
\end{minipage}
\end{enumerate}
\vspace{-0.3\baselineskip}
\begin{small}
\textit{คำตอบ: ข}
\end{small}

\item เมื่อเพิ่มความเข้มของแสง (โดยความถี่คงที่) ผลที่เกิดขึ้นกับปรากฏการณ์โฟโตอิเล็กทริกคืออะไร?
\begin{enumerate}[label=\alph*)., itemsep=0.3\baselineskip, topsep=0.3\baselineskip]
\item ก\ 
\begin{minipage}[t]{0.9\linewidth}
\raggedright พลังงานจลน์สูงสุดของอิเล็กตรออนเพิ่มขึ้น
\end{minipage}
\item ข\ 
\begin{minipage}[t]{0.9\linewidth}
\raggedright จำนวนอิเล็กตรออนที่หลุดต่อวินาทีเพิ่มขึ้น
\end{minipage}
\item ค\ 
\begin{minipage}[t]{0.9\linewidth}
\raggedright ไม่มีผลอะไร
\end{minipage}
\item ง\ 
\begin{minipage}[t]{0.9\linewidth}
\raggedright อิเล็กตรออนหลุดเร็วขึ้น
\end{minipage}
\end{enumerate}
\vspace{-0.3\baselineskip}
\begin{small}
\textit{คำตอบ: ข}
\end{small}

\item แสง UV ทำให้เกิดปรากฏการณ์โฟโตอิเล็กทริกได้ง่ายกว่าแสงที่ตามองเห็น เพราะอะไร?
\begin{enumerate}[label=\alph*)., itemsep=0.3\baselineskip, topsep=0.3\baselineskip]
\item ก\ 
\begin{minipage}[t]{0.9\linewidth}
\raggedright UV มีความเข้มมากกว่า
\end{minipage}
\item ข\ 
\begin{minipage}[t]{0.9\linewidth}
\raggedright UV มีความถี่สูงกว่า จึงมีพลังงานโฟตอนมากกว่า
\end{minipage}
\item ค\ 
\begin{minipage}[t]{0.9\linewidth}
\raggedright UV มีความยาวคลื่นมากกว่า
\end{minipage}
\item ง\ 
\begin{minipage}[t]{0.9\linewidth}
\raggedright UV สะท้อนได้ดีกว่า
\end{minipage}
\end{enumerate}
\vspace{-0.3\baselineskip}
\begin{small}
\textit{คำตอบ: ข}
\end{small}

\item ถ้าโมเมนตัมของอิเล็กตรออนวัดได้แม่นยำมาก (Δp ≈ 0) ข้อใดจะเป็นจริง?
\begin{enumerate}[label=\alph*)., itemsep=0.3\baselineskip, topsep=0.3\baselineskip]
\item ก\ 
\begin{minipage}[t]{0.9\linewidth}
\raggedright ตำแหน่งของอิเล็กตรออนจะแน่นอนแม่นยำ
\end{minipage}
\item ข\ 
\begin{minipage}[t]{0.9\linewidth}
\raggedright ตำแหน่งของอิเล็กตรออนจะไม่แน่นอนมาก (Δx มาก)
\end{minipage}
\item ค\ 
\begin{minipage}[t]{0.9\linewidth}
\raggedright อิเล็กตรออนจะหยุดนิ่ง
\end{minipage}
\item ง\ 
\begin{minipage}[t]{0.9\linewidth}
\raggedright หลักความไม่แน่นอนจะไม่มีผล
\end{minipage}
\end{enumerate}
\vspace{-0.3\baselineskip}
\begin{small}
\textit{คำตอบ: ข}
\end{small}

\item ในกลศาสตร์ควอนตัม สถานะของอิเล็กตรออนในอะตอมถูกอธิบายด้วยอะไร?
\begin{enumerate}[label=\alph*)., itemsep=0.3\baselineskip, topsep=0.3\baselineskip]
\item ก\ 
\begin{minipage}[t]{0.9\linewidth}
\raggedright วงโคจรที่แน่นอน
\end{minipage}
\item ข\ 
\begin{minipage}[t]{0.9\linewidth}
\raggedright ฟังก์ชันคลื่นหรือ orbital
\end{minipage}
\item ค\ 
\begin{minipage}[t]{0.9\linewidth}
\raggedright ตำแหน่งที่แน่นอน
\end{minipage}
\item ง\ 
\begin{minipage}[t]{0.9\linewidth}
\raggedright การสั่นภายในนิวเคลียร์
\end{minipage}
\end{enumerate}
\vspace{-0.3\baselineskip}
\begin{small}
\textit{คำตอบ: ข}
\end{small}

\item อิเล็กตรออนและโปรตอนเคลื่อนที่ด้วยความเร็วเท่ากัน ข้อใดถูกต้องเกี่ยวกับความยาวคลื่น de Broglie ของทั้งสอง?
\begin{enumerate}[label=\alph*)., itemsep=0.3\baselineskip, topsep=0.3\baselineskip]
\item ก\ 
\begin{minipage}[t]{0.9\linewidth}
\raggedright เท่ากัน
\end{minipage}
\item ข\ 
\begin{minipage}[t]{0.9\linewidth}
\raggedright อิเล็กตรออนมีความยาวคลื่นมากกว่า
\end{minipage}
\item ค\ 
\begin{minipage}[t]{0.9\linewidth}
\raggedright โปรตอนมีความยาวคลื่นมากกว่า
\end{minipage}
\item ง\ 
\begin{minipage}[t]{0.9\linewidth}
\raggedright ไม่สามารถเปรียบเทียบได้
\end{minipage}
\end{enumerate}
\vspace{-0.3\baselineskip}
\begin{small}
\textit{คำตอบ: ข}
\end{small}

\item โลหะโซเดียมมีค่างาน 2.3 eV จะมีความถี่ขีดเริ่มเท่าใด?
\begin{enumerate}[label=\alph*)., itemsep=0.3\baselineskip, topsep=0.3\baselineskip]
\item ก\ 
\begin{minipage}[t]{0.9\linewidth}
\raggedright 5.5 × 10¹⁴ Hz
\end{minipage}
\item ข\ 
\begin{minipage}[t]{0.9\linewidth}
\raggedright 5.5 × 10¹³ Hz
\end{minipage}
\item ค\ 
\begin{minipage}[t]{0.9\linewidth}
\raggedright 2.3 × 10¹⁵ Hz
\end{minipage}
\item ง\ 
\begin{minipage}[t]{0.9\linewidth}
\raggedright 2.3 × 10¹⁴ Hz
\end{minipage}
\end{enumerate}
\vspace{-0.3\baselineskip}
\begin{small}
\textit{คำตอบ: ข}
\end{small}

\item ถ้าโฟตอนมีพลังงาน 3.0 eV ความยาวคลื่นของมันเป็นเท่าใด?
\begin{enumerate}[label=\alph*)., itemsep=0.3\baselineskip, topsep=0.3\baselineskip]
\item ก\ 
\begin{minipage}[t]{0.9\linewidth}
\raggedright 207 nm
\end{minipage}
\item ข\ 
\begin{minipage}[t]{0.9\linewidth}
\raggedright 310 nm
\end{minipage}
\item ค\ 
\begin{minipage}[t]{0.9\linewidth}
\raggedright 414 nm
\end{minipage}
\item ง\ 
\begin{minipage}[t]{0.9\linewidth}
\raggedright 518 nm
\end{minipage}
\end{enumerate}
\vspace{-0.3\baselineskip}
\begin{small}
\textit{คำตอบ: ค}
\end{small}

\item นิวตรออนถูกกักไว้ในนิวเคลียร์ขนาด \textasciitilde{}10⁻¹⁴ m ความไม่แน่นอนของโมเมนตัมจะมีค่าอย่างน้อยประมาณเท่าใด?
\begin{enumerate}[label=\alph*)., itemsep=0.3\baselineskip, topsep=0.3\baselineskip]
\item ก\ 
\begin{minipage}[t]{0.9\linewidth}
\raggedright 5 × 10⁻²⁰ kg·m/s
\end{minipage}
\item ข\ 
\begin{minipage}[t]{0.9\linewidth}
\raggedright 5 × 10⁻²² kg·m/s
\end{minipage}
\item ค\ 
\begin{minipage}[t]{0.9\linewidth}
\raggedright 5 × 10⁻²⁴ kg·m/s
\end{minipage}
\item ง\ 
\begin{minipage}[t]{0.9\linewidth}
\raggedright 5 × 10⁻²⁶ kg·m/s
\end{minipage}
\end{enumerate}
\vspace{-0.3\baselineskip}
\begin{small}
\textit{คำตอบ: ข}
\end{small}

\item ข้อใดต่อไปนี้ไม่ใช่ปรากฏการณ์ที่ต้องอธิบายด้วยกลศาสตร์ควอนตัม?
\begin{enumerate}[label=\alph*)., itemsep=0.3\baselineskip, topsep=0.3\baselineskip]
\item ก\ 
\begin{minipage}[t]{0.9\linewidth}
\raggedright ปรากฏการณ์โฟโตอิเล็กทริก
\end{minipage}
\item ข\ 
\begin{minipage}[t]{0.9\linewidth}
\raggedright การเลี้ยวเบนของอิเล็กตรออน
\end{minipage}
\item ค\ 
\begin{minipage}[t]{0.9\linewidth}
\raggedright การเคลื่อนที่ของดาวเคราะห์รอบดวงอาทิตย์
\end{minipage}
\item ง\ 
\begin{minipage}[t]{0.9\linewidth}
\raggedright สเปกตรัมเส้นของอะตอมไฮโดรเจน
\end{minipage}
\end{enumerate}
\vspace{-0.3\baselineskip}
\begin{small}
\textit{คำตอบ: ค}
\end{small}

\item ถ้าอิเล็กตรออนถูกเร่งด้วยความต่างศักย์ 100 V ความยาวคลื่น de Broglie จะเปลี่ยนจากเท่าใดเป็นเท่าใด?
\begin{enumerate}[label=\alph*)., itemsep=0.3\baselineskip, topsep=0.3\baselineskip]
\item ก\ 
\begin{minipage}[t]{0.9\linewidth}
\raggedright เพิ่มขึ้น
\end{minipage}
\item ข\ 
\begin{minipage}[t]{0.9\linewidth}
\raggedright ลดลง
\end{minipage}
\item ค\ 
\begin{minipage}[t]{0.9\linewidth}
\raggedright คงที่
\end{minipage}
\item ง\ 
\begin{minipage}[t]{0.9\linewidth}
\raggedright ขึ้นอยู่กับทิศทาง
\end{minipage}
\end{enumerate}
\vspace{-0.3\baselineskip}
\begin{small}
\textit{คำตอบ: ข}
\end{small}

\item เมื่อใช้แสงเลเซอร์สีแดง (λ = 650 nm, พลังงานต่ำ) ยิงที่โลหะทังสเตน (φ = 4.5 eV) จะเกิดอะไรขึ้น?
\begin{enumerate}[label=\alph*)., itemsep=0.3\baselineskip, topsep=0.3\baselineskip]
\item ก\ 
\begin{minipage}[t]{0.9\linewidth}
\raggedright อิเล็กตรออนหลุดทันที
\end{minipage}
\item ข\ 
\begin{minipage}[t]{0.9\linewidth}
\raggedright อิเล็กตรออนไม่หลุดเพราะพลังงานโฟตอนน้อยกว่าค่างาน
\end{minipage}
\item ค\ 
\begin{minipage}[t]{0.9\linewidth}
\raggedright อิเล็กตรออนหลุดช้า
\end{minipage}
\item ง\ 
\begin{minipage}[t]{0.9\linewidth}
\raggedright โลหะเปล่งแสง
\end{minipage}
\end{enumerate}
\vspace{-0.3\baselineskip}
\begin{small}
\textit{คำตอบ: ข}
\end{small}

\item ความถี่คี่ต่ำสุดที่ทำให้อิเล็กตรออนหลุดจากโลหะที่มีค่างาน 3.0 eV ต้องมีความยาวคลื่นไม่เกินเท่าใด?
\begin{enumerate}[label=\alph*)., itemsep=0.3\baselineskip, topsep=0.3\baselineskip]
\item ก\ 
\begin{minipage}[t]{0.9\linewidth}
\raggedright 207 nm
\end{minipage}
\item ข\ 
\begin{minipage}[t]{0.9\linewidth}
\raggedright 310 nm
\end{minipage}
\item ค\ 
\begin{minipage}[t]{0.9\linewidth}
\raggedright 414 nm
\end{minipage}
\item ง\ 
\begin{minipage}[t]{0.9\linewidth}
\raggedright 518 nm
\end{minipage}
\end{enumerate}
\vspace{-0.3\baselineskip}
\begin{small}
\textit{คำตอบ: ค}
\end{small}

\item ข้อใดไม่สามารถใช้หลักความไม่แน่นอนอธิบายได้?
\begin{enumerate}[label=\alph*)., itemsep=0.3\baselineskip, topsep=0.3\baselineskip]
\item ก\ 
\begin{minipage}[t]{0.9\linewidth}
\raggedright ทำไมอิเล็กตรออนไม่ตกลงนิวเคลียร์
\end{minipage}
\item ข\ 
\begin{minipage}[t]{0.9\linewidth}
\raggedright การกระจายตัวของอิเล็กตรออนในช่องแคบคู่
\end{minipage}
\item ค\ 
\begin{minipage}[t]{0.9\linewidth}
\raggedright สเปกตรัมเส้นของไฮโดรเจน
\end{minipage}
\item ง\ 
\begin{minipage}[t]{0.9\linewidth}
\raggedright ความหนาแน่นของนิวตรออนในนิวเคลียร์
\end{minipage}
\end{enumerate}
\vspace{-0.3\baselineskip}
\begin{small}
\textit{คำตอบ: ค}
\end{small}

\item ในการทดลองช่องแคบคู่ (double-slit experiment) กับอิเล็กตรออน เมื่อยิงทีละตัวจะเกิดอะไรขึ้น?
\begin{enumerate}[label=\alph*)., itemsep=0.3\baselineskip, topsep=0.3\baselineskip]
\item ก\ 
\begin{minipage}[t]{0.9\linewidth}
\raggedright ไม่เกิดรูปแบบเลี้ยวเบน
\end{minipage}
\item ข\ 
\begin{minipage}[t]{0.9\linewidth}
\raggedright เกิดจุดบนฉากทีละจุด แต่เมื่อรวมจุดหลายจุดจะเห็นรูปแบบเลี้ยวเบน
\end{minipage}
\item ค\ 
\begin{minipage}[t]{0.9\linewidth}
\raggedright อิเล็กตรออนผ่านได้เพียงช่องเดียวเสมอ
\end{minipage}
\item ง\ 
\begin{minipage}[t]{0.9\linewidth}
\raggedright อิเล็กตรออนแต่ละตัวเลือกช่องทางใดทางหนึ่ง
\end{minipage}
\end{enumerate}
\vspace{-0.3\baselineskip}
\begin{small}
\textit{คำตอบ: ข}
\end{small}

\item อิเล็กตรออนในกล่องกว้าง 1 nm มีความยาวคลื่น de Broglie ประมาณเท่าใด?
\begin{enumerate}[label=\alph*)., itemsep=0.3\baselineskip, topsep=0.3\baselineskip]
\item ก\ 
\begin{minipage}[t]{0.9\linewidth}
\raggedright 0.5 nm
\end{minipage}
\item ข\ 
\begin{minipage}[t]{0.9\linewidth}
\raggedright 1 nm
\end{minipage}
\item ค\ 
\begin{minipage}[t]{0.9\linewidth}
\raggedright 2 nm
\end{minipage}
\item ง\ 
\begin{minipage}[t]{0.9\linewidth}
\raggedright 5 nm
\end{minipage}
\end{enumerate}
\vspace{-0.3\baselineskip}
\begin{small}
\textit{คำตอบ: ข}
\end{small}

\end{enumerate}
\end{document}